\documentclass[11pt,a4paper]{article}

\usepackage[utf8]{inputenc}
\usepackage[T1]{fontenc}
\usepackage{amsmath,amssymb,amsfonts}
\usepackage{physics}
\usepackage{graphicx}
\usepackage{booktabs}
\usepackage{array}
\usepackage{siunitx}
\usepackage{hyperref}
\usepackage[margin=1in]{geometry}
\usepackage{xcolor}
\usepackage{listings}
\usepackage{enumitem}
\usepackage{fancyhdr}
\usepackage{titlesec}

\hypersetup{
    colorlinks=true,
    linkcolor=blue!60!black,
    urlcolor=blue!60!black,
    citecolor=blue!60!black,
    pdftitle={PY-BEMCS Physics Reference and User Manual},
    pdfauthor={Dr. Bharat Singh Rawat, Nick Magrin}
}

\pagestyle{fancy}
\fancyhf{}
\fancyhead[L]{PY-BEMCS Reference Manual}
\fancyhead[R]{\thepage}
\fancyfoot[C]{\small Python Beam Extraction \& Monte Carlo Simulator}

\titleformat{\section}{\Large\bfseries\color{blue!50!black}}{\thesection}{1em}{}
\titleformat{\subsection}{\large\bfseries\color{blue!40!black}}{\thesubsection}{1em}{}
\titleformat{\subsubsection}{\normalsize\bfseries\color{blue!30!black}}{\thesubsubsection}{1em}{}

\newcommand{\PYBEMCS}{PY-BEMCS}

\lstset{
    basicstyle=\small\ttfamily,
    backgroundcolor=\color{gray!10},
    frame=single,
    breaklines=true,
    captionpos=b
}

\begin{document}

% ---------------------- TITLE ----------------------
\begin{titlepage}
\centering
\vspace*{1.5cm}

{\Huge\bfseries Python Beam Extraction \&\\[0.3em] Monte Carlo Simulator \par}
\vspace{0.4cm}
{\Large\bfseries (\PYBEMCS)\par}

\vspace{1.0cm}
{\large\textit{Comprehensive Physics Reference, Numerical Controls, and User Manual}\par}

\vspace{1.8cm}
{\large\textbf{Original Author:}\par}
{\large Dr.~Bharat Singh Rawat\par}
{\normalsize \href{mailto:bharat.bharat22@gmail.com}{bharat.bharat22@gmail.com}\par}

\vspace{0.8cm}
{\large\textbf{Version 2.0 Extended Architecture \& Physics Modifications:}\par}
{\large Nick Magrin\par}
{\normalsize \href{mailto:nick.magrin@gmail.com}{nick.magrin@gmail.com}\par}

\vspace{1.2cm}
{\normalsize \textbf{Version 2.0 (Extended Modular Release)}\par}
\vspace{0.2cm}
{\normalsize \today\par}

\vfill
\hrule
\vspace{0.5em}
{\small Open-Source 2D3V Electrostatic Particle-in-Cell, Plume MCC, Thermo-Mechanical Grid Erosion, Numerical Verification Controls, \& Multi-Core Benchmarking Suite}
\end{titlepage}

% ---------------------- ABSTRACT ----------------------
\begin{center}
{\Large\bfseries Abstract}
\end{center}
\vspace{0.5em}

\PYBEMCS{} is an open-source 2D3V (two spatial, three velocity components) electrostatic Particle-in-Cell (PIC) and Monte Carlo Collision (MCC) simulation suite engineered for ion thruster optics, neutral beam injectors (NBI), and advanced plasma extraction systems. This document provides a complete mathematical, physical, and numerical specification of the simulator, together with rigorous stability criteria, physics validation checks, and comprehensive user guidelines for the graphical interface, headless batch solvers, and automated multi-core benchmarking scripts.

Key capabilities detailed in this manual include:
\begin{itemize}[nosep]
    \item Self-consistent Bohm extraction from upstream plasma sheaths with space-charge Poisson solves and Boltzmann electron screening.
    \item Multi-zone non-uniform finite-volume mesh with automatic geometry partitioning and adaptive axial coarsening factors ($1.0\times, 1.5\times, 2.0\times, 4.0\times$).
    \item Adaptive Picard iteration with dynamic residual convergence ($\Delta V_{\max} \le \SI{50}{\milli\volt}$) and sticky Newton--Raphson fallback solver.
    \item Automatic numerical physics validation controls (Debye length resolution, plasma frequency limits, CFL velocity bounds, and dynamic domain scaling).
    \item 2D3V relativistic Boris particle advancement with dual Taichi GPU/CPU accelerated and pure-NumPy fallback kernels.
    \item Automated parametric sweep framework with real-time multi-scenario switching and multi-metric steady-state convergence criteria.
    \item Charge-Exchange (CEX) collision dynamics using Roy's neutral plume model or user-imported experimental cross-section datasets with log-log cubic spline smoothing.
    \item In-situ sputter erosion modeling (Matsunami yield formulation) coupled with live grid mesh boundary deletion.
    \item Coupled 2D thermal conduction, Stefan--Boltzmann radiative cooling, and Euler--Bernoulli cantilever thermo-mechanical grid aperture deflection.
    \item Modular JSON configuration architecture (\texttt{config.json}) supporting arbitrary $N$-grid configurations, discharge chamber plasma dynamics, and live parameter reloading.
    \item Automated multi-core parallel benchmarking suite for CEX erosion rates, Electron Backstreaming (EBS) limits, primary beam impingement / grid scraping, and perveance-divergence curves.
    \item Built-in asynchronous standalone application packager (\texttt{PyInstallerWorker}) embedded directly in the GUI.
\end{itemize}

\vspace{1em}
\hrule
\vspace{1em}

\tableofcontents
\newpage

% =====================================================
\section{Computational Framework and Architecture}
% =====================================================

\subsection{Coordinate System and Multi-Zone Non-Uniform Discretisation}

\PYBEMCS{} supports two distinct spatial formulations configured via \texttt{"coordinate\_system"}:
\begin{enumerate}[nosep]
    \item \textbf{Cartesian 2D Slab (\texttt{"cartesian"}, Default)}: Discretises a planar 2D slice $(x, y)$ with an assumed invariant unit transverse depth $\Delta z = \SI{1}{\milli\meter}$ (slit optics).
    \item \textbf{Axisymmetric 2D3V Cylindrical (\texttt{"axisymmetric"})}: Discretises a true 3D circular beamlet in cylindrical coordinates $(z, r)$, mapped as $x \equiv z$ (axial beam extraction coordinate) and $y \equiv r \ge 0$ (radial coordinate from the beamlet centerline).
\end{enumerate}

To optimize computational throughput while strictly resolving local physical scales (such as the sub-millimeter upstream Debye length $\lambda_{\mathrm{D}}$ and fine aperture grid chamfers), \PYBEMCS{} implements a \textbf{multi-zone non-uniform axial mesh} coupled with a uniform transverse spacing:
\begin{itemize}[nosep]
    \item Transverse spacing remains strictly uniform: $\Delta y = \text{const}$ across all $j \in [0, n_y - 1]$, where $n_y = \lfloor L_y / \Delta y \rfloor + 1$.
    \item Axial spacing $\Delta x(x)$ adapts across automatically detected physical zones:
    \begin{enumerate}[nosep]
        \item \textbf{Presheath Zone} ($x \in [0, x_{\mathrm{screen\_start}}]$): Baseline fine resolution $f_{\mathrm{pre}} = 1.0\times$ ($\Delta x = \Delta x_0 \le \lambda_{\mathrm{D}}$) to accurately capture Bohm meniscus formation.
        \item \textbf{Optics Grid Barrels} ($x \in [x_k, x_k + t_k]$): Moderate resolution $f_{\mathrm{barrel}} = 1.0\text{--}1.5\times$ ensuring $\ge 15$ cells across grid thickness $t_k$ for high-fidelity boundary tracking and sputter erosion.
        \item \textbf{Inter-Grid Acceleration Gaps} ($x \in [x_k + t_k, x_{k+1}]$): $f_{\mathrm{gap}} = 1.0\text{--}2.0\times$ resolving strong electrostatic acceleration fields while maintaining electric field gradient jumps below $10\%$.
        \item \textbf{Near and Far Plume Zones} ($x \in [x_{\mathrm{exit}}, L_x]$): Coarsened resolution $f_{\mathrm{plume}} = 2.0\text{--}4.0\times$, exploiting the dilute quasi-neutral state of the downstream beam to achieve a $25\text{--}35\%$ reduction in total grid nodes.
    \end{enumerate}
\end{itemize}

For each axial node $i$, the non-uniform control volume width $h_i^x$ is defined by:
\begin{equation}
    h_i^x = \frac{\Delta x_{i-1/2} + \Delta x_{i+1/2}}{2}, \qquad \Delta x_{i+1/2} = x_{i+1} - x_i.
\end{equation}
In Cartesian mode, the cell volume is given by $\Omega_{i,j} = h_i^x \Delta y \Delta z$. In Axisymmetric mode, each nodal cell represents a complete annular ring of revolution around the beamlet axis $r = 0$:
\begin{equation}
    \Omega_{i,j} = V_{\mathrm{rad}, j} \cdot h_i^x, \qquad
    V_{\mathrm{rad}, j} = 
    \begin{cases}
        \pi \left(\frac{\Delta r}{2}\right)^2, & j = 0 \text{ (on-axis node)}, \\
        2\pi r_j \Delta r, & 1 \le j \le n_y - 2 \text{ (interior annuli)}, \\
        \pi \left[r_{n_y-1}^2 - \left(r_{n_y-1} - \frac{\Delta r}{2}\right)^2\right], & j = n_y - 1 \text{ (outer boundary)}.
    \end{cases}
\end{equation}
This rigorous volumetric weighting ensures exact machine-precision charge conservation ($|\sum Q_p - \iiint \rho\, \dd{V}| / \sum Q_p < 10^{-8}$) in both Cartesian and cylindrical coordinates.

Every macroparticle stores three-dimensional velocity components $\vb{v} = (v_x, v_y, v_z)$. In axisymmetric mode, $(v_x, v_y, v_z)$ correspond to the axial velocity $v_z$, radial velocity $v_r$, and azimuthal velocity $v_\theta$.

\subsection{Dual-Kernel Execution Engine (Taichi GPU + CPU Fallback)}

To combine peak GPU performance on workstation hardware with maximum portability across containerized or headless computing clusters, \PYBEMCS{} implements a dual-kernel architecture:
\begin{enumerate}[nosep]
    \item \textbf{Taichi Backend (GPU/CPU)}: When Taichi is initialized, particle pushes, bilinear charge deposition, and explicit thermal conduction compile into parallel GPU kernels (CUDA, Metal, or Vulkan) running in 32-bit floating-point precision ($\texttt{ti.f32}$).
    \item \textbf{Vectorised CPU Fallback Kernels}: If Taichi is not present (or when running inside frozen executable packages), the physics engine automatically activates vectorised CPU routines (\texttt{accumulate\_rho\_cpu}, \texttt{push\_particles\_boris\_cpu}) using optimized NumPy array operations.
\end{enumerate}

% =====================================================
\section{Rigorous Numerical Physics Controls and Stability Verification}
% =====================================================

To prevent unphysical numerical heating, artificial grid-aliasing instabilities, or phase-space skipping, the physics engine enforces strict automated validation checks during domain generation (\texttt{build\_domain}):

\subsection{Debye Length Resolution Criterion}
To prevent unphysical finite-grid self-heating, the spatial cell dimensions must resolve the upstream plasma Debye length:
\begin{equation}
    \lambda_{\mathrm{D}} = \sqrt{\frac{\varepsilon_0 T_{\mathrm{e,up}}}{e \cdot (0.61 n_0)}},
\end{equation}
where the pre-sheath factor $0.61 \approx \exp(-1/2)$ accounts for the density drop at the sheath edge. The engine strictly requires:
\begin{equation}
    \Delta x \le \lambda_{\mathrm{D}} \quad \text{and} \quad \Delta y \le \lambda_{\mathrm{D}}.
\end{equation}
If the cell spacing exceeds $\lambda_{\mathrm{D}}$, a descriptive error is raised, prompting the user to refine $\Delta x, \Delta y$ or adjust the plasma density.

\subsection{Plasma Frequency and Time-Step Integration Limits}
The time step $\Delta t$ must resolve both the ion plasma frequency $\omega_{\mathrm{pi}}$ and the scaled electron plasma frequency $\omega_{\mathrm{pe}}$:
\begin{equation}
    \omega_{\mathrm{pi}} = \sqrt{\frac{n_0 e^2}{m_{\mathrm{ion}} \varepsilon_0}}, \qquad
    \omega_{\mathrm{pe}} = \sqrt{\frac{n_0 e^2}{m_{\mathrm{e}} \varepsilon_0}}.
\end{equation}
The maximum permissible time step is bounded by:
\begin{equation}
    \Delta t \le \frac{2\pi}{\omega_{\mathrm{pi}}}, \qquad \Delta t_{\mathrm{e}} \le \frac{2\pi}{\omega_{\mathrm{pe}}}.
\end{equation}
If $\Delta t$ violates this limit, domain initialization halts to avoid numerical divergence.

\subsection{CFL and Velocity Resolution Bound}
Particles must not travel more than one mesh cell during a single time increment. The maximum expected particle velocity at the sheath edge includes the Bohm drift and 4 standard deviations of the thermal distribution:
\begin{equation}
    v_{\max} = v_{\mathrm{B}} + 4 v_{\mathrm{th,i}} = \sqrt{\frac{Z e T_{\mathrm{e}}}{m_{\mathrm{ion}}}} + 4\sqrt{\frac{Z e T_{\mathrm{i}}}{m_{\mathrm{ion}}}}.
\end{equation}
The engine enforces the Courant--Friedrichs--Lewy (CFL) velocity resolution bound:
\begin{equation}
    \frac{\Delta x}{\Delta t} \ge v_{\max} \quad \text{and} \quad \frac{\Delta y}{\Delta t} \ge v_{\max}.
\end{equation}

\subsection{Dynamic Domain Sizing and Symmetry Boundary Conditions}
The axial computational domain length $L_x$ is dynamically calculated based on the total physical grid stack:
\begin{equation}
    L_x = \SI{1.5}{\milli\meter} + \sum_{k=1}^{N} \left(t_k + g_k\right) + \SI{3.0}{\milli\meter},
\end{equation}
ensuring adequate space upstream for meniscus formation and downstream for plume expansion.

The transverse domain height $L_y$ and boundary conditions adapt to the selected geometry layout and coordinate system:
\begin{itemize}[nosep]
    \item \textbf{\texttt{half\_hole}} (Default): Half-pitch symmetry representation with specular reflection at the transverse boundary $y = 0$ (and $y = L_y$ when uncoupled), where $L_y = 0.5 r_{\mathrm{s}} + 0.30 \cdot \text{pitch}$ or the hexagonal Wigner--Seitz unit-cell radius $R_{\mathrm{WS}} \approx 0.525 \cdot \text{pitch}$.
    \item \textbf{\texttt{one\_hole}}: Full single-aperture extraction domain ($L_y = 3.0 r_{\mathrm{s}}$) with periodic boundary conditions along $y$ (Cartesian 2D slab only).
    \item \textbf{\texttt{two\_holes}}: Dual-aperture inter-hole coupling domain ($L_y = 3.0 r_{\mathrm{s}} + \text{pitch}$) with periodic boundaries along $y$ (Cartesian 2D slab only).
\end{itemize}

\subsubsection{Geometric and Physical Incompatibility of Multi-Hole Modes in Axisymmetric ($r\text{--}z$) Optics}
When the solver operates in \textbf{Axisymmetric ($r\text{--}z$)} cylindrical mode, the coordinate system is governed by rotation of the 2D computational domain around the singular centerline $r = 0$ ($y = 0$). Consequently, \PYBEMCS{} strictly constrains the geometry option to \textbf{\texttt{half\_hole}} and blocks \texttt{one\_hole} and \texttt{two\_holes} for fundamental physical and geometrical reasons:
\begin{enumerate}[nosep]
    \item \textbf{Rotational Symmetry Generation}: In cylindrical coordinates, revolving any 2D planar point $(z, r)$ by $2\pi$ generates a 3D circle of radius $r$ centered on the axis. If an aperture were centered at $r_0 > 0$ (as in \texttt{one\_hole} or \texttt{two\_holes}), revolving this aperture around $r = 0$ would physically generate a \textbf{hollow toroidal annular slit (a donut-shaped channel)} with solid metal on the thruster centerline, rather than a discrete circular ion optics aperture.
    \item \textbf{Breakdown of Azimuthal Invariance for Adjacent Apertures}: Simulating two adjacent circular beamlets side-by-side fundamentally breaks azimuthal symmetry ($\partial / \partial \theta \ne 0$). In an $r\text{--}z$ code, two apertures at $r_1$ and $r_2$ revolve into two concentric nested toroidal rings. Modeling multiple discrete circular beamlets requires a full 3D Cartesian solver (such as the \texttt{Cpp3D} module).
    \item \textbf{Singular Axis at $r = 0$ vs. Periodic Boundaries}: The centerline $r = 0$ represents the singular axis of rotational symmetry, where $E_r(0, z) = 0$ and crossing ions experience specular radial velocity reversal ($r \leftarrow -r \implies r \ge 0,\; v_r \leftarrow -v_r$). In contrast, \texttt{one\_hole} and \texttt{two\_holes} enforce periodic boundary conditions ($V(0) = V(L_y)$), which would unphysically map particles exiting at the outer radius $r = L_y$ to the centerline $r = 0$, coupling annular cells of vastly different volumes ($V_{\mathrm{cell}} \propto r$) and violating Gauss's law and charge conservation.
\end{enumerate}

% =====================================================
\subsubsection{Transverse Injection Span and Sheath Expansion Controls}
In Section 1 (Geometry / Grid Optics) of the user interface, two parameters govern the transverse domain sizing and upstream injection region:
\begin{itemize}[nosep]
    \item \textbf{Hole-Boundary Margin (\texttt{hole\_boundary\_dist\_mm})}: Enforces a minimum radial/transverse solid webbing clearance between the outermost grid aperture wall and the upper computational domain boundary at  = L_y$.
    \item \textbf{Injection Y-Length (\texttt{inj\_y\_length\_mm})}: Configures the transverse or radial span $[0, y_{\mathrm{inj}}]$ over which upstream ions (and RF electrons) are introduced at the injection boundary  = 0$. By default, {\mathrm{inj}}$ matches the screen grid aperture radius {\mathrm{s}}$ (with an \texttt{Auto} button to restore synchronization). However, users can increase {\mathrm{inj}}$ up to the full domain height $. This models physical Debye sheath expansion ($\lambda_{\mathrm{D}} \propto n_0^{-1/2}$) over the solid screen grid webbing at low plasma densities, funneling peripheral ions into the aperture at steep incidence angles to accurately capture low-density crossover divergence.
\end{itemize}

\section{Particle-in-Cell Electrostatic Algorithm}
% =====================================================

\subsection{Self-Consistent Poisson--Boltzmann Solver}

The electrostatic potential $\phi(\vb{x})$ is governed by Poisson's equation including macroparticle ion density and thermalized Boltzmann electrons:
\begin{equation}
    \nabla^{2}\phi(\vb{x}) \;=\; -\,\frac{1}{\varepsilon_{0}}
        \left[ q_{\mathrm{i}}\,n_{\mathrm{i}}(\vb{x}) - e\,n_{\mathrm{e}}(\vb{x}) \right],
    \label{eq:poisson}
\end{equation}
where $q_{\mathrm{i}} = Z e$ is the ion charge, $n_{\mathrm{i}}$ is the deposited ion macroparticle number density, and the electron density $n_{\mathrm{e}}$ follows the non-linear Boltzmann relation:
\begin{equation}
    n_{\mathrm{e}}(\vb{x}) \;=\; n_{0}\,
    \exp\!\left[\frac{\min(\phi(\vb{x}),\,\phi_{\mathrm{pl}}) - \phi_{\mathrm{pl}}}{T_{\mathrm{e}}}\right],
    \label{eq:boltzmann}
\end{equation}
with $T_{\mathrm{e}}$ in electron-volts ($\si{\eV}$), $n_0$ the upstream plasma density, and $\phi_{\mathrm{pl}}$ the reference plasma potential established at the upstream boundary ($x = 0$). The clipping operator $\min(\phi, \phi_{\mathrm{pl}})$ prevents non-physical exponential density blow-ups in regions where local potential excursions exceed the upstream plasma potential.

In \textbf{Cartesian 2D slab mode}, Equation~\eqref{eq:poisson} is discretised on the non-uniform mesh using a finite-volume formulation. Multiplying by $\Delta y^2$ yields the scaled 5-point stencil:
\begin{equation}
    c_W \phi_{i-1,j} + c_E \phi_{i+1,j} + \phi_{i,j-1} + \phi_{i,j+1} - (c_W + c_E + 2.0)\phi_{i,j}
    = -\frac{\Delta y^2 \rho_{i,j}}{\varepsilon_{0}},
\end{equation}
where $c_W = \frac{\Delta y^2}{h_i^x \Delta x_{i-1/2}}$ and $c_E = \frac{\Delta y^2}{h_i^x \Delta x_{i+1/2}}$.

In \textbf{Axisymmetric ($r\text{--}z$) cylindrical mode}, the Laplacian operator expands to:
\begin{equation}
    \nabla^2 \phi = \frac{\partial^2 \phi}{\partial z^2} + \frac{1}{r}\frac{\partial}{\partial r}\left(r \frac{\partial \phi}{\partial r}\right) = \frac{\partial^2 \phi}{\partial z^2} + \frac{\partial^2 \phi}{\partial r^2} + \frac{1}{r}\frac{\partial \phi}{\partial r}.
\end{equation}
Multiplying by $\Delta r^2$ yields the interior 5-point stencil for nodes $j \ge 1$:
\begin{equation}
    c_W \phi_{i-1,j} + c_E \phi_{i+1,j} + \left(1 - \frac{1}{2j}\right)\phi_{i,j-1} + \left(1 + \frac{1}{2j}\right)\phi_{i,j+1} - (c_W + c_E + 2.0)\phi_{i,j}
    = -\frac{\Delta r^2 \rho_{i,j}}{\varepsilon_{0}}.
\end{equation}
At the centerline axis of revolution $r = 0$ ($j = 0$), the geometric singularity is evaluated via L'H\^opital's rule:
\begin{equation}
    \lim_{r \to 0} \frac{1}{r}\frac{\partial \phi}{\partial r} = \frac{\partial^2 \phi}{\partial r^2} \implies \nabla^2 \phi\Big|_{r=0} = \frac{\partial^2 \phi}{\partial z^2} + 2\frac{\partial^2 \phi}{\partial r^2}.
\end{equation}
Applying radial symmetry $\phi(z, -\Delta r) = \phi(z, \Delta r)$ yields the exact discrete on-axis finite-volume row:
\begin{equation}
    c_W \phi_{i-1,0} + c_E \phi_{i+1,0} + 4.0 \phi_{i,1} - (c_W + c_E + 4.0)\phi_{i,0} = -\frac{\Delta r^2 \rho_{i,0}}{\varepsilon_0},
\end{equation}
with on-axis electric field boundary condition $E_r(z, 0) = 0$. Scaling both Cartesian and axisymmetric operators by $\Delta y^2 = \Delta r^2$ preserves an identical right-hand side source vector $\vb{b} = -\frac{\Delta r^2}{\varepsilon_0}\boldsymbol{\rho}$, allowing the unified adaptive Picard, Anderson Acceleration, and Newton--Raphson solver pipelines to operate identically across both coordinate systems. The resulting sparse operator $\mathbf{A} \boldsymbol{\phi} = \mathbf{b}$ is pre-factorised using SuperLU sparse LU decomposition (or GPU-accelerated CuPy sparse LU).

\subsubsection{Adaptive Picard Iteration and Sticky Newton--Raphson Fallback}

To guarantee numerical stability and eliminate unphysical electric field spikes, the non-linear Boltzmann electron response is solved adaptively:
\begin{enumerate}[nosep]
    \item \textbf{Adaptive Picard Under-Relaxation}: At iteration $k$, the right-hand side space charge vector $\mathbf{b}^{(k)}$ incorporates the non-linear electron fluid density $\rho_e(\boldsymbol{\phi}^{(k)})$. The intermediate solution $\boldsymbol{\phi}^* = \mathbf{A}^{-1} \mathbf{b}^{(k)}$ is relaxed via:
    \begin{equation}
        \boldsymbol{\phi}^{(k+1)} = (1 - \omega) \boldsymbol{\phi}^{(k)} + \omega \boldsymbol{\phi}^*, \qquad \omega = 0.2.
    \end{equation}
    Iterations proceed dynamically until the maximum potential residual satisfies:
    \begin{equation}
        \|\Delta \boldsymbol{\phi}\|_{\infty} = \max_{i,j} \left|\phi_{i,j}^{(k+1)} - \phi_{i,j}^{(k)}\right| \le \SI{50}{\milli\volt}.
    \end{equation}
    \item \textbf{Sticky Newton--Raphson Fallback}: If Picard iteration stagnates (residual decreases by less than $1\%$ over 5 iterations) or exceeds 35 iterations, the engine automatically switches to a linearized Newton--Raphson solver:
    \begin{equation}
        \mathbf{J}(\boldsymbol{\phi}^{(k)}) \, \delta \boldsymbol{\phi} = -\mathbf{F}(\boldsymbol{\phi}^{(k)}), \qquad \boldsymbol{\phi}^{(k+1)} = \boldsymbol{\phi}^{(k)} + \delta \boldsymbol{\phi},
    \end{equation}
    where the Jacobian diagonal contains the Boltzmann derivative $J_{m,m} = A_{m,m} + \frac{\Delta y^2 e n_e(\phi_m)}{\varepsilon_0 T_e}$. Once engaged, the solver remains in Newton mode ('sticky') until steady-state convergence is re-established, achieving quadratic convergence without divergence risk.
\end{enumerate}

\subsection{Bilinear Charge Deposition}

Macroparticle charge is mapped to the four surrounding mesh nodes using bilinear interpolation weights. For a particle located at $(x_p, y_p)$ with grid indices $i = \lfloor x_p/\Delta x \rfloor$, $j = \lfloor y_p/\Delta y \rfloor$:
\begin{align}
    h_x &= \frac{x_p - i\Delta x}{\Delta x}, \qquad
    h_y = \frac{y_p - j\Delta y}{\Delta y},\\
    \rho_{i,j}     &\mathrel{+}= \frac{q_{\mathrm{i}} w}{V_{\mathrm{cell}}}\,(1-h_x)(1-h_y),\\
    \rho_{i+1,j}   &\mathrel{+}= \frac{q_{\mathrm{i}} w}{V_{\mathrm{cell}}}\,h_x(1-h_y),\\
    \rho_{i,j+1}   &\mathrel{+}= \frac{q_{\mathrm{i}} w}{V_{\mathrm{cell}}}\,(1-h_x)h_y,\\
    \rho_{i+1,j+1} &\mathrel{+}= \frac{q_{\mathrm{i}} w}{V_{\mathrm{cell}}}\,h_x h_y,
\end{align}
where $V_{\mathrm{cell}} = \Delta x \Delta y \Delta z$ is the cell volume with assumed unit transverse thickness $\Delta z = \SI{1}{\milli\meter}$, and $w$ is the statistical macroparticle weighting factor:
\begin{equation}
    w = \max\!\left(\frac{n_0 V_{\mathrm{cell}}}{N_{\mathrm{ppc}}},\; 10^{3}\right),
\end{equation}
with default nominal target $N_{\mathrm{ppc}} = 40$ particles per cell.

\subsection{Electric Field Computation and Boris Pusher}

Electric field components are evaluated at mesh nodes via second-order central differencing:
\begin{equation}
    E_{x,i,j} = -\frac{\phi_{i+1,j} - \phi_{i-1,j}}{2\Delta x}, \qquad
    E_{y,i,j} = -\frac{\phi_{i,j+1} - \phi_{i,j-1}}{2\Delta y},
\end{equation}
and interpolated to particle positions. Particle trajectories are advanced using the explicit leap-frog Boris method~\cite{birdsall}:
\begin{align}
    \vb{v}^{-} &= \vb{v}^{n-1/2} + \frac{q\,\vb{E}}{2m}\,\Delta t,\\
    \vb{t} &= \frac{q\,\vb{B}}{2m}\,\Delta t, \qquad
    \vb{s} = \frac{2\,\vb{t}}{1 + |\vb{t}|^2},\\
    \vb{v}' &= \vb{v}^{-} + \vb{v}^{-} \times \vb{t},\\
    \vb{v}^{+} &= \vb{v}^{-} + \vb{v}' \times \vb{s},\\
    \vb{v}^{n+1/2} &= \vb{v}^{+} + \frac{q\,\vb{E}}{2m}\,\Delta t,\\
    \vb{x}^{n+1} &= \vb{x}^{n} + \vb{v}^{n+1/2}\,\Delta t.
\end{align}
This scheme guarantees exact conservation of phase-space volume and energy conservation in static magnetic fields.

\subsubsection{Cylindrical Birdsall \& Langdon Coordinate Rotation}
In \textbf{Axisymmetric ($r\text{--}z$) mode}, standard leap-frog advancement in cylindrical coordinates would introduce an artificial $1/r$ centrifugal acceleration singularity ($v_\theta^2 / r$). To eliminate coordinate singularities while strictly preserving canonical angular momentum $L_z = m r v_\theta$ and transverse kinetic energy $E_\perp = \frac{1}{2}m(v_r^2 + v_\theta^2)$, \PYBEMCS{} implements the \textbf{Birdsall \& Langdon local-Cartesian rotation}~\cite{birdsall} in both CPU and Taichi GPU kernels:
\begin{align}
    Y &= r^n + v_r^{n+1/2}\,\Delta t, \qquad Z = v_\theta^{n+1/2}\,\Delta t,\\
    r^{n+1} &= \sqrt{Y^2 + Z^2},\\
    \cos\alpha &= \frac{Y}{r^{n+1}}, \qquad \sin\alpha = \frac{Z}{r^{n+1}},\\
    v_r^{n+1/2} &\leftarrow v_r^{n+1/2} \cos\alpha + v_\theta^{n+1/2} \sin\alpha,\\
    v_\theta^{n+1/2} &\leftarrow -v_r^{n+1/2} \sin\alpha + v_\theta^{n+1/2} \cos\alpha.
\end{align}
If an ion passes through the axis $r = 0$, its radius reflects ($r \ge 0$) and its radial velocity specularly reverses ($v_r \leftarrow -v_r$), naturally maintaining beamlet symmetry.

% =====================================================
\section{Plasma Injection and Bohm Optics}
% =====================================================

\subsection{Bohm Sheath Extraction Criterion}

Ions enter the computational domain through the upstream plasma boundary ($x = 0$) with an initial axial drift satisfying the Bohm sheath velocity criterion:
\begin{equation}
    v_{\mathrm{B}} = \sqrt{\frac{Z e T_{\mathrm{e}}}{m_{\mathrm{ion}}}}, \qquad
    m_{\mathrm{ion}} = M_{\mathrm{amu}} \cdot u,
    \label{eq:bohm}
\end{equation}
where $M_{\mathrm{amu}}$ is the ion mass in atomic mass units and $u = \SI{1.66054e-27}{\kilo\gram}$.

The injected ion current across the open sheath area $A_{\mathrm{inj}}$ is:
\begin{equation}
    I_{\mathrm{ion}} = 0.61\, Z e n_{0} v_{\mathrm{B}} A_{\mathrm{inj}},
\end{equation}
where the prefactor $0.61 \approx \exp(-1/2)$ accounts for the quasi-neutral pre-sheath density drop. In Cartesian mode, $A_{\mathrm{inj}} = \Delta y_{\mathrm{inj}} \cdot \Delta z$ (with unit depth $\Delta z = \SI{1}{\milli\meter}$) and initial transverse coordinates are sampled uniformly. In Axisymmetric mode, $A_{\mathrm{inj}} = \pi r_{\mathrm{inj}}^2$, and macroparticle radial injection positions are sampled according to:
\begin{equation}
    r = r_{\mathrm{inj}} \sqrt{U}, \qquad U \sim \mathcal{U}(0, 1),
\end{equation}
guaranteeing a strictly uniform number density per unit cross-sectional area across all concentric cylindrical annuli. The number of macroparticles injected at each timestep $\Delta t$ is:
\begin{equation}
    N_{\mathrm{inj}} = \frac{I_{\mathrm{ion}} \Delta t}{Z e w} + \xi_{\mathrm{stochastic}},
\end{equation}
with stochastic remainder accumulation ensuring exact long-term current conservation.

\subsection{Thermal Velocity Sampling}

Thermal spread at the sheath boundary is sampled from a 3D Maxwellian distribution at ion temperature $T_{\mathrm{i}}$:
\begin{equation}
    v_x = v_{\mathrm{B}} + \zeta_x v_{\mathrm{th,i}}, \qquad
    v_y = \zeta_y v_{\mathrm{th,i}}, \qquad
    v_z = \zeta_z v_{\mathrm{th,i}},
\end{equation}
where $\zeta_{x,y,z} \sim \mathcal{N}(0, 1)$ are standard normal random variables, and $v_{\mathrm{th,i}} = \sqrt{Z e T_{\mathrm{i}} / m_{\mathrm{ion}}}$.

% =====================================================
\section{Multiphysics Grid Models}
% =====================================================

\subsection{Arbitrary $N$-Grid Geometry Specification}

The extraction optics consist of an arbitrary sequence of $N$ conducting grids defined by:
\begin{itemize}[nosep]
    \item DC potential $V_k$ (\si{\volt})
    \item Thickness $t_k$ (\si{\milli\meter})
    \item Inter-grid axial gap to subsequent grid $g_k$ (\si{\milli\meter})
    \item Aperture radius $r_k$ (\si{\milli\meter})
    \item Upstream chamfer angle $\theta_k$ (\si{\degree})
\end{itemize}

Solid boundary masks $\mathcal{M}_k(x,y)$ are dynamically synthesised and rasterised onto the finite-difference mesh.

\subsection{Cantilever Thermo-Mechanical Grid Deformation}

Under intense ion impingement and radiative heating, individual grid membranes experience substantial temperature rises $\Delta T_k = T_k - T_{\infty}$. Each grid is modeled as an annular plate segment undergoing Euler--Bernoulli cantilever bending clamped at the outer boundary ($y = L_y$) and free at the aperture edge ($y = r_k$):
\begin{equation}
    \delta_{\max,k} = \frac{\alpha_k \Delta T_k L_k^2}{2 t_k}, \qquad L_k = L_y - r_k,
\end{equation}
where $\alpha_k$ is the linear thermal expansion coefficient ($\si{\per\kelvin}$). The transverse displacement profile:
\begin{equation}
    \delta(y) = \delta_{\max,k} \left(\frac{L_y - y}{L_k}\right)^2,
\end{equation}
is dynamically added as an axial coordinate distortion to the grid boundaries. Whenever cumulative structural deformation exceeds $\SI{5}{\micro\meter}$, the Poisson sparse matrix operator is automatically re-factored.

\subsection{2D Thermal Diffusion and Radiative Dissipation}

Thermal energy deposition from impacting ions is balanced by 2D conduction within the grid solid mask and Stefan--Boltzmann radiative cooling to space:
\begin{equation}
    \rho_k c_{p,k} \frac{\partial T}{\partial t} = k_k \left(\frac{\partial^2 T}{\partial x^2} + \frac{\partial^2 T}{\partial y^2}\right) + \dot{q}_{\mathrm{dep}} - \frac{\varepsilon_k \sigma_{\mathrm{SB}}}{\Delta z}\left(T^4 - T_{\infty}^4\right),
\end{equation}
where $k_k$ is thermal conductivity, $c_{p,k}$ is specific heat capacity, $\varepsilon_k$ is surface emissivity, $\sigma_{\mathrm{SB}} = \SI{5.67e-8}{\watt\per\meter\squared\per\kelvin\tothe{4}}$, and $T_{\infty} = \SI{300}{\kelvin}$.

% =====================================================
\section{Collisions and Sputter Erosion}
% =====================================================

\subsection{Charge-Exchange (CEX) Dynamics and Plume Neutral Model}

Neutral gas density distribution throughout the grid gap and downstream region follows Roy's analytical plume expansion model~\cite{roy}:
\begin{equation}
    n_{\mathrm{n}}(x,y) = n_{\mathrm{n},0} \, a \left[1 - \frac{1}{\sqrt{1 + (r_T / R)^2}}\right] \cos\theta,
\end{equation}
with $R = \sqrt{y^2 + (x + r_T)^2}$, $\theta = \arctan(y / (x + r_T))$, and normalisation factor $a = (1 - 1/\sqrt{2})^{-1}$.

The null-collision Monte Carlo collision (MCC) probability for an ion traversing distance $\Delta s = |\vb{v}|\Delta t$ is:
\begin{equation}
    P_{\mathrm{CEX}} = 1 - \exp\left[-n_{\mathrm{n}}(x,y) \, \sigma_{\mathrm{CEX}}(E) \, |\vb{v}| \, \Delta t\right].
\end{equation}
The analytical xenon CEX cross-section formulation is:
\begin{equation}
    \sigma_{\mathrm{CEX}}(g) = \left[-0.8821\ln(g) + 15.1262\right]^2 \times 10^{-20} \;\si{\meter\squared},
\end{equation}
where $g = |\vb{v}|$ is relative collision speed. If custom cross-section CSV tables are imported, this formula is dynamically superseded by cubic spline interpolation in $\log_{10}(E)$--$\log_{10}(\sigma)$ space.

\subsection{Sputter Yield and Dynamic Cell Removal}

Erosion of grid boundaries under ion bombardment is governed by the empirical Matsunami sputtering yield relation~\cite{matsunami}:
\begin{equation}
    Y(E) = \min\!\left[Y_0 \left(E_{\mathrm{eV}} - E_{\mathrm{th}}\right)^{1.5},\; 1.0\right],
\end{equation}
where $Y_0$ is the sputter yield coefficient and $E_{\mathrm{th}}$ is the sputtering threshold energy ($\si{\eV}$).

Cumulative damage per cell is accumulated as:
\begin{equation}
    D_{i,j} \mathrel{+}= Y(E) \cdot \kappa_{\mathrm{acc}},
\end{equation}
where $\kappa_{\mathrm{acc}}$ is the numerical time-acceleration factor. When $D_{i,j} \ge D_{\mathrm{fail}}$, the grid cell is removed from the solid mask, exposing new interior nodes and updating the field solver in real time.

% =====================================================
\section{Modular Configuration Architecture (\texttt{config.json})}
% =====================================================

The simulator parameters are managed through a unified, hierarchical JSON structure (\texttt{config.json}). A representative configuration is shown below:

\begin{lstlisting}[language=json,caption={Standard \texttt{config.json} structure}]
{
  "beam_species": {
    "mass_amu": 131.293,
    "charge_state": 1
  },
  "grid_material": {
    "preset": "Molybdenum"
  },
  "simulation": {
    "coordinate_system": "axisymmetric",
    "n0_plasma": 7.73e17,
    "Te_up": 8.5,
    "Ti": 0.034,
    "Tn": 400.0,
    "n0": 5.27e18,
    "Accel": 0.5,
    "Thresh": 10000.0,
    "sim_mode": "Both",
    "geometry": "half_hole"
  },
  "discharge_chamber": {
    "radius_m": 0.02,
    "length_m": 0.10,
    "n_cusp": 4,
    "pitch_mm": 3.0
  },
  "grids": [
    { "V": 1300.0, "t": 1.2, "gap": 0.9, "r": 1.3, "cham": 20.0 },
    { "V": -150.0, "t": 0.6, "gap": 0.5, "r": 0.8, "cham": 0.0 },
    { "V": 50.0,   "t": 0.5, "gap": 4.0, "r": 1.2, "cham": 0.0 }
  ],
  "cross_sections": {
    "cx_file": "",
    "see_file": "",
    "custom_file": "",
    "spline_smoothing": 0.0
  },
  "terminal_output": {
    "grid_temperatures": true,
    "beam_divergence": true,
    "saddle_point_potential": true,
    "mean_particle_energy": true,
    "iteration_time": true
  }
}
\end{lstlisting}

% =====================================================
\section{Graphical User Interface and Features}
% =====================================================

\subsection{Menu Bar and Advanced Controls}

\begin{itemize}
    \item \textbf{Settings $\rightarrow$ Advanced Parameters...}: Configures neutralizer injection coordinates $(x_{\mathrm{n}}, r_{\mathrm{n}})$, plasma offset potential, and electron mass scaling ratio $\mathcal{R} = m_{\mathrm{ion}}/m_{\mathrm{e}}$.
    \item \textbf{Settings $\rightarrow$ Open Config JSON...}: Interactive file browser to load external JSON configuration profiles.
    \item \textbf{Settings $\rightarrow$ Reload config.json}: Hot-reloads configuration files from disk without restarting the application.
    \item \textbf{Settings $\rightarrow$ Build .exe...}: Compiles the current Python source code into a standalone single-file binary executable using an asynchronous background thread with live milestone tracking.
    \item \textbf{Beam $\rightarrow$ Ion Species...}: Configures beam ion atomic mass (amu) and ionization charge state ($+1, +2, \dots$) with predefined presets (Xe, Kr, Ar, N$_2$, O$_2$, H$_2$, He, Hg, Cs).
    \item \textbf{Beam $\rightarrow$ Cross-Section Manager...}: Import external experimental CSV tables for CEX, SEE yield, or custom reaction channels with interactive log-log cubic spline fitting.
    \item \textbf{Materials $\rightarrow$ Grid Material...}: Selects material presets (Molybdenum, Steel SS316, Titanium, Graphite) or custom thermophysical properties ($k, \rho, c_p, \varepsilon, \alpha, E, Y_0, E_{\mathrm{th}}$).
\end{itemize}

\subsection{Comprehensive Diagnostic Windows and Telemetry Suite}

The simulator provides an extensive suite of standalone, interactive diagnostic windows equipped with native Matplotlib navigation toolbars (Pan, Zoom, Save) and live data streaming:
\begin{itemize}[nosep]
    \item \textbf{Beam Divergence Angle ($\theta_{95}$ \& Midline)}: Evaluated from the 95\textsuperscript{th} percentile (or 90\textsuperscript{th}, Direct RMS, RMS Twiss, FWHM) trajectory angles at configurable axial evaluation planes downstream of the optics. In Cartesian mode:
    \begin{equation}
        \theta_{95} = Q_{0.95}\!\left[\,\left|\arctan\left(\frac{v_y}{v_x}\right)\right|\,\right].
    \end{equation}
    In Axisymmetric mode, the true 3D beamlet cone divergence is evaluated incorporating the azimuthal velocity $v_\theta = v_z$:
    \begin{equation}
        v_\perp = \sqrt{v_r^2 + v_\theta^2}, \qquad \theta_{\mathrm{3D}} = \arctan2(v_\perp, v_x) \cdot \frac{180^\circ}{\pi}, \qquad \theta_{95} = Q_{0.95}\left[\theta_{\mathrm{3D}}\right].
    \end{equation}
    \item \textbf{2D Plasma Number Density Visualizer (\texttt{density\_window.py})}: Full 2D logarithmic color contours of ion and electron density maps coupled with 1D radial profiles cut along two independent downstream axial planes.
    \item \textbf{Steady-State Convergence Monitor (\texttt{steady\_state\_window.py})}: Continuous rolling window variance analysis tracking key parameters ($\theta_{95}$, $V_{\mathrm{saddle}}$, $I_{\mathrm{beam}}$, $E_{\mathrm{kin}}$). Steady state is confirmed when rolling relative fluctuations drop below $\le 1.0\%$ over a configurable evaluation duration.
    \item \textbf{Electron Backstreaming (EBS) Saddle Potential}: Continuous monitoring of centerline potential in the accelerator aperture to verify backstreaming suppression ($V_{\mathrm{sp}} < -5\,\si{\volt}$).
    \item \textbf{Mesh Discretization \& Sub-Step Inspector (\texttt{mesh\_window.py})}: Visualizes axial zone boundaries, cell size transitions $\Delta x(x)$, and the real-time distribution of Boris sub-steps required by high-velocity particles.
    \item \textbf{Energy \& Charge Budget Conservation (\texttt{energy\_window.py}, \texttt{charge\_window.py})}: Tracks instantaneous kinetic, potential, and field energy conservation, as well as macroparticle vs.\ Boltzmann electron charge neutrality.
    \item \textbf{Ion Energy Distribution Function (\texttt{iedf\_window.py})}: Live spectral histogram tracking particle kinetic energy distributions exiting the extraction system.
    \item \textbf{Cross-Section Manager (\texttt{cs\_viewer\_window.py})}: Log--log visualizer and cubic-spline interpolation inspector for imported reaction cross-section tables.
    \item \textbf{Asynchronous Data Exporters}: Background progress dialogs for exporting time-series simulation telemetry to CSV, particle phase-space archives, and animated GIF recordings.
\end{itemize}

% =====================================================
\section{Multi-Core Benchmarking Suite (\texttt{benchmarks/})}
% =====================================================

The \texttt{benchmarks/} directory contains dedicated multi-core physics sweep scripts utilising Python's \texttt{multiprocessing} engine to execute parametric studies across independent CPU worker processes simultaneously:

\begin{table}[h!]
\centering
\small
\begin{tabular}{lp{6cm}p{4.5cm}}
\toprule
\textbf{Benchmark Script} & \textbf{Physics Objective} & \textbf{Generated Output} \\
\midrule
\texttt{benchmark\_cex.py} & Sputter erosion rate vs background neutral density ($n_{\mathrm{n}} = 10^{18}$--$5\times 10^{19}\,\si{\per\meter\cubed}$) across multiple grid gap configurations. & \texttt{cex\_benchmark\_multicore.png} \\
\addlinespace
\texttt{benchmark\_ebs.py} & Centerline saddle-point potential barrier vs accelerator grid voltage ($V_{\mathrm{a}} = -400$--$0\,\si{\volt}$) evaluating the EBS safety boundary. & \texttt{ebs\_benchmark\_multicore.png} \\
\addlinespace
\texttt{benchmark\_impingement.py} & Direct primary beam ion grid interception and scraping percentage vs perveance ($P = I_{\mathrm{ion}} / V_{\mathrm{tot}}^{1.5}$). & \texttt{impingement\_ratio\_benchmark.png} \\
\addlinespace
\texttt{benchmark\_perveance.py} & Beam divergence angle vs perveance over 60 plasma density increments across multiple grid gap spacings. & \texttt{divergence\_vs\_perveance\_benchmark\_multicore.png} \\
\addlinespace
\texttt{benchmark\_perveance\_Vs\_Sweep.py} & Beam divergence response as a function of screen grid extraction potential ($V_{\mathrm{s}} = 600$--$2000\,\si{\volt}$). & \texttt{divergence\_vs\_voltage\_benchmark.png} \\
\bottomrule
\end{tabular}
\caption{Summary of multi-core benchmarking scripts.}
\end{table}

% =====================================================
\section{Standalone Executable Compilation Architecture}
% =====================================================

The embedded executable builder (\texttt{Settings $\rightarrow$ Build .exe...}) runs PyInstaller asynchronously via the \texttt{PyInstallerWorker(QThread)} class. The compilation process proceeds without blocking the Qt event loop:
\begin{enumerate}[nosep]
    \item Spawns a dedicated background worker thread and builds inside a temporary directory (\texttt{\_exe\_build\_tmp/}) to keep the repository root clean.
    \item Intercepts real-time standard output and maps build milestones (analysis, dependency collection, PYZ archive creation, binary packaging) to a 0--100\% progress dialog.
    \item Automatically copies the resulting single-file executable and accompanying \texttt{config.json} into the user-selected destination directory.
    \item Cleans up intermediate build artifacts upon completion or cancellation.
\end{enumerate}

% =====================================================
\section{Summary of Execution Modes}
% =====================================================

\subsection{Interactive Graphical Run}
\begin{lstlisting}[language=bash]
python Python/main.py
\end{lstlisting}

\subsection{Headless Parameter Sweep Execution}
\begin{lstlisting}[language=bash]
python Python/run_simulation_from_config.py
\end{lstlisting}

\subsection{Parallel Multi-Core Benchmarks}
\begin{lstlisting}[language=bash]
python Python/benchmarks/benchmark_perveance.py
python Python/benchmarks/benchmark_cex.py
python Python/benchmarks/benchmark_ebs.py
\end{lstlisting}

% =====================================================
\begin{thebibliography}{9}
\bibitem{birdsall}
C.~K. Birdsall and A.~B. Langdon,
\textit{Plasma Physics via Computer Simulation},
Institute of Physics Publishing, Bristol, 1991.

\bibitem{roy}
R.~I.~S. Roy, \textit{Numerical Simulation of Ion Thruster Plume Backflow}, Ph.D. Thesis, Massachusetts Institute of Technology, Cambridge, MA, 1995.

\bibitem{matsunami}
N.~Matsunami \textit{et al.},
``Energy dependence of the yields of ion-induced sputtering of monatomic solids,'' \textit{Atomic Data and Nuclear Data Tables}, vol.~31, no.~1, pp.~1--80, 1984.
\end{thebibliography}

\end{document}
